% !TEX root = ../../main.tex
% C8.4 - Search quality evaluation. 5 Oct 2026: 26 queries (files/wildtrack_evaluation_queries_v2.json), current
% configuration, backend/var/evaluation/wildtrack-v2-conf010.json (B3 in files/part1-improvement-plan.md).
\section{Search Quality Evaluation}
\label{sec:8.4}

\paragraph{Method.} The pipeline was run with the current configuration (YOLO11n at a confidence threshold of 0.1, ByteTrack, RaSa, every frame) on the labelled frames of the seven WILDTRACK cameras, 400 frames per camera, which produced a gallery of 1,526 tracks; each track was given, for evaluation only, the majority person among its matched labels. The evaluation set has 26 queries, each in three forms: a cropped image, an English description and a set of attributes. Six of them were written for the first version of the evaluation; twenty identities were added on 5 October, chosen among people seen on all seven cameras with at least 30 labelled boxes, with descriptions and attributes written from three crops on different cameras using only the attribute vocabulary of the application; two identities walking next to a person dressed alike were left out. The observation used as the query image is removed from the gallery, and every other labelled observation of the same person is a correct result, 5 to 33 per query. \emph{Recall@$k$} is the fraction of queries with at least one correct result in the top $k$, with $k \in \{4, 8, 12, 16\}$, and \emph{MRR} is the mean reciprocal rank of the first correct result; text and attribute queries are tokenised with RaSa's own tokenizer (Section~\ref{sec:3.5.4}). With 26 queries one query is 3.8 points of recall, so Table~\ref{tab:recall} gives the 95\% Wilson interval of Recall@8. An earlier run of the first six queries at a confidence threshold of 0.25, on a gallery of 1,552 tracks, gave the same picture: Recall@8 of 1.0 by image and 0 by text and by attributes.

\begin{table}[H]
\centering
\caption{Search quality results at the current configuration (26 queries, 1,526 tracks, confidence threshold 0.1)}
\label{tab:recall}
\begin{tabularx}{\textwidth}{|L|c|c|c|c|c|c|}
\hline
\textbf{Form} & \textbf{R@4} & \textbf{R@8} & \textbf{R@12} & \textbf{R@16} & \textbf{MRR} & \textbf{95\% interval of R@8} \\
\hline
By image & 0.846 & 0.885 & 0.885 & 0.885 & 0.671 & 0.71 to 0.96 \\
\hline
By text & 0 & 0 & 0.038 & 0.038 & 0.021 & 0 to 0.13 \\
\hline
By attributes & 0 & 0 & 0 & 0.038 & 0.015 & 0 to 0.13 \\
\hline
\end{tabularx}
\end{table}

\paragraph{Person detection and tracking.} At the confidence threshold of 0.1, detection on the 2,800 labelled frames reaches a precision of 0.43 and a recall of 0.64 at an IoU of 0.5 (27,194 correct boxes, 35,602 extra boxes, 15,527 labelled people missed); the earlier run at 0.25 gave 0.58 and 0.54. The lower threshold misses fewer people but adds false detections, mostly removed by tracking (Section~\ref{sec:3.2.5}): of the 1,526 tracks, 491 match no labelled person, and 825 of the 1,639 labelled people are covered by at least one track.

\paragraph{Discussion.} Search by image performs well: 23 of the 26 queries have a correct result in the top 8, the first correct result is at rank 1 for 14 queries, and the three failures are at ranks 99, 109 and 553. Search by text and by attributes, in contrast, bring almost nothing into the top 16: one text query (at rank 12) and one attribute query (at rank 15). The first correct result lies at a median rank of 73 by text and 111 by attributes in the gallery of 1,526 tracks; with 5 to 33 correct tracks per query, a random order would place it at a median rank of about 150, so the text encoder separates very little on this data. Two causes are established below and in Section~\ref{sec:3.5.4}: the application searches with RaSa's contrastive vectors alone, while the model's accuracy relies on its cross-modal re-ranker, and the domain gap between CUHK-PEDES, on which RaSa was trained, and the person images cropped automatically from the full frames of WILDTRACK. Search by image is therefore the main form, while text and attribute search play a supporting role that needs careful visual assessment. This weakness of text search is the starting point of the research in Part~II (Chapter~\ref{chapter:r-overview}).

\paragraph{Pipeline check on the training domain.} To tell an implementation fault from the domain gap, the adapters that serve queries (image bytes through the image encoder, sentence through the text encoder, inner-product ranking) were run on 600 images of the CUHK-PEDES training split, the domain of the checkpoint, with their captions; the published test split is not redistributed, and the copy used carries the captions of a person on each of that person's images, so retrieval is scored per identity. Three things follow. Image-to-image retrieval of the same person is perfect (Recall@1 of 1.000 on 109 queries), so the image side is right. Text-to-image retrieval by the vectors alone, as the application searches, reaches Recall@1 of 0.035 and Recall@10 of 0.177, against 0.002 and 0.017 for a random order: well above chance, far below the published 76.5\%. Re-ranking the top 32 candidates with RaSa's cross-modal matching head, as the authors' evaluation code does over the top 128, raises Recall@1 to 0.650 on 40 of these queries over a 150-image gallery (contrastive vectors alone: 0.050), at roughly 50 milliseconds per candidate on the CPU. The published accuracy of RaSa therefore belongs to its re-ranker, which the application keeps disabled (Section~\ref{sec:3.5.4}); the contrastive vectors that the vector search compares are only its candidate generator. This, not an implementation fault, is the first reason text search is weak, before any domain gap. The check also found one difference from the authors' code: queries were tokenised with the standard BERT tokenizer, which closes a sentence with [SEP], while RaSa trains and evaluates without it; using RaSa's own tokenizer raised Recall@10 on this sample from 0.137 to 0.177 and is now the default, with the image vectors and the index unchanged.

\paragraph{Qualitative check on the demonstration data.} On the demonstration data (Section~\ref{sec:7.2}), five query images were tried with the top 8 results and assessed visually: three queries had 8/8 correct results, one 7/8 and one 4/8; the correct results lay on 4 to 5 different cameras. This is only a qualitative check on a dataset different from the evaluation set and cannot be compared with the Recall values of Table~\ref{tab:recall}.
